\chapter{Production Checklist}
\label{ch:bf-production-checklist}

This checklist is the operating definition of BayesFilter production readiness.
It is intentionally stricter than a successful smoke test.

\section{Target and Source}
\label{sec:bf-production-target-source}

Before HMC:
\begin{enumerate}[label=(\roman*)]
  \item the scalar log target is defined in writing;
  \item every formula has provenance in source notes, local code/tests, or
    reviewed literature;
  \item parameter transforms and Jacobian terms are included;
  \item invalid structural regions return documented finite values;
  \item approximate filters are labeled as target-defining approximations or
    corrected surrogates.
\end{enumerate}

\section{Derivative and Filter}
\label{sec:bf-production-derivative-filter}

The filter backend must provide:
\begin{enumerate}[label=(\roman*)]
  \item shape, dtype, mask, and static-compilation contracts;
  \item finite value and finite gradient checks;
  \item derivative parity against finite differences or autodiff on controlled
    cases;
  \item backend parity where multiple linear algebra forms exist;
  \item factor, solve, condition-number, and spectral-gap diagnostics;
  \item an explicit custom-gradient policy when raw autodiff is not the final
    production path.
\end{enumerate}

\section{Sampler and JIT}
\label{sec:bf-production-sampler-jit}

The sampler must report:
\begin{enumerate}[label=(\roman*)]
  \item whether the complete value-and-gradient path is compiled;
  \item eager/compiled parity on small targets;
  \item mass-matrix source and conditioning;
  \item acceptance rate, divergences, E-BFMI or equivalent energy diagnostics,
    R-hat, ESS, MCSE/SD, and runtime;
  \item saved JSON and chain artifacts for medium and strict runs;
  \item sidecar capture artifacts only for bounded debugging, not for every
    production transition.
\end{enumerate}

\section{Claim Discipline}
\label{sec:bf-production-claim-discipline}

Use the narrowest honest label:
\begin{description}
  \item[value-valid] the likelihood value is finite and reference-checked;
  \item[gradient-valid] the derivative matches the target value and passes
    parity checks;
  \item[sampler-usable] finite HMC diagnostics justify the next experiment;
  \item[converged] strict multi-chain diagnostics pass and artifacts are saved;
  \item[production-ready] all preceding gates pass under the intended model,
    data, hardware, and compilation policy.
\end{description}
Anything less should stay in the reset memo as a gap, not disappear into prose.

\section{Fixed-Center Geometry Checklist}
\label{sec:bf-production-fixed-center-geometry}

Before a fixed-center precision enters HMC:
\begin{enumerate}[label=(\roman*)]
  \item use \texttt{diagnostic\_center}, not \texttt{map}, unless a separate
    MAP gate passed;
  \item preserve raw dense precision, eigenvalues, inertia, and projection
    burden before SPD regularization;
  \item require at least two fit replicates and caller-reviewed pairwise caps;
  \item keep training, selection holdout, and audit clouds disjoint;
  \item verify covariance/precision orientation and coordinate scaling;
  \item reject unidentified two-factor fits and record anchor/rank diagnostics;
  \item follow the fallback order and never manufacture identity automatically;
  \item construct a \texttt{PrecomputedMassArtifact} only after every gate;
  \item run a separate exact Metropolis-corrected HMC mechanics canary;
  \item exclude adaptation and tuning draws from retained inference.
\end{enumerate}

For CPU/XLA cloud evaluation, record the explicit route, core-count source,
worker count or override, $B=1$, XLA flag, child PIDs, deterministic ordering,
and semantic progress events.  Never convert a GPU failure into automatic CPU
fallback.  GPU execution must verify incremental memory growth before framework
initialization.

Passing this checklist only nominates an exact HMC canary.  It does not
establish posterior convergence or production readiness.  Target-only or
child-function XLA success does not establish full-chain XLA.
